% --- Archon physics formalization source begin ---
% archon:physics
% archon:covers PhyXMiniProblems/problem_phyx_mini_0560.lean
% archon:source-report reports/phyx_mini/problem_phyx_mini_0560.source.json

\chapter{Physics problem phyx\_mini\_0560}
\label{ch:PhyXMiniProblems_problem_phyx_mini_0560}

\paragraph{Problem source.}
\#\# Physical scenario

Suppose that a particle with mass \textbackslash{}( m \textbackslash{}) and energy \textbackslash{}( E \textbackslash{}) is moving toward the origin of a system \textbackslash{}( S \textbackslash{}) such that its velocity \textbackslash{}( \textbackslash{}mathbf\{u\} \textbackslash{}) makes an angle \textbackslash{}( \textbackslash{}alpha \textbackslash{}) with the \textbackslash{}( y \textbackslash{}) axis as shown in figure.

\#\# Question

determine the energy \textbackslash{}( E' \textbackslash{}) of the particle measured by an observer in \textbackslash{}( S' \textbackslash{}), which moves relative to \textbackslash{}( S \textbackslash{}) so that the particle moves along the \textbackslash{}( y' \textbackslash{}) axis.

\#\# Answer choices

- A: A: \textbackslash{}( E\textbackslash{}sqrt\{1 + (u\textasciicircum{}2 / c\textasciicircum{}2)\textbackslash{}sin\textasciicircum{}2 \textbackslash{}alpha\} \textbackslash{})

- B: B: \textbackslash{}( E\textbackslash{}sqrt\{2 - (u\textasciicircum{}2 / c\textasciicircum{}2)\textbackslash{}sin\textasciicircum{}2 \textbackslash{}alpha\} \textbackslash{})

- C: C: \textbackslash{}( E\textbackslash{}sqrt\{2 + (u\textasciicircum{}2 / c\textasciicircum{}2)\textbackslash{}sin\textasciicircum{}2 \textbackslash{}alpha\} \textbackslash{})

- D: D: \textbackslash{}( E\textbackslash{}sqrt\{1 - (u\textasciicircum{}2 / c\textasciicircum{}2)\textbackslash{}sin\textasciicircum{}2 \textbackslash{}alpha\} \textbackslash{})

\#\# Figure caption (auxiliary; use the image as primary evidence)

The image depicts a two-dimensional Cartesian coordinate system with x and y axes. There is an arrow representing a vector, labeled **u**, that originates from the origin and extends into the first quadrant. The vector **u** is marked with a point labeled **m** near its tip.

There is a blue arrow indicating the direction of the vector, and it forms an angle labeled **α** with the positive x-axis. The text **S** is present in the background, suggesting an area or region in the coordinate space. The vector is shown with a solid line and extends diagonally.

The axes are labeled with **x** and **y**. The relationship between the vector **u** and the angle **α** as well as the coordinate axes are visually represented.

\#\# Dataset metadata

Relativity; Spatial Relation Reasoning, Multi-Formula Reasoning

\paragraph{Recorded answer/context.}
D: D: \textbackslash{}( E\textbackslash{}sqrt\{1 - (u\textasciicircum{}2 / c\textasciicircum{}2)\textbackslash{}sin\textasciicircum{}2 \textbackslash{}alpha\} \textbackslash{})

\paragraph{Figure/image path.}
/root/proposal\_for\_physic/science-mango/phyx\_mini\_run/phyx\_data/test\_image/560.png

\paragraph{Formalization target.}
create a compiling Lean file with sorry bodies at `PhyXMiniProblems/problem\_phyx\_mini\_0560.lean`. The Lean declarations must preserve the physical quantities, dimensions or dimensional roles, figure labels, governing-law hypotheses, and final relation expressed by this problem.
Use Mathlib/Physlib names found through LeanExplore where available. If a physics API is missing, introduce faithful local abstractions rather than scalar placeholder aliases.

\begin{theorem}[Physics formalization target]
\label{thm:physics:phyx_mini_0560:target}
\lean{PhyXMiniProblems.ProblemPhyXMini0560.recorded_answer_D_matches}
\uses{def:physics:phyx-mini-0560:phyxminiproblems-problemphyxmini0560-transverseboostsetup, def:physics:phyx-mini-0560:phyxminiproblems-problemphyxmini0560-matchesproblemandprimaryfigure, def:physics:phyx-mini-0560:phyxminiproblems-problemphyxmini0560-hasphysicalinputparameters, def:physics:phyx-mini-0560:phyxminiproblems-problemphyxmini0560-observermakesparticlemovealongyprime, def:physics:phyx-mini-0560:phyxminiproblems-problemphyxmini0560-satisfiesspecialrelativisticparticlekinematics, def:physics:phyx-mini-0560:phyxminiproblems-problemphyxmini0560-recordeddatasetanswer, def:physics:phyx-mini-0560:phyxminiproblems-problemphyxmini0560-matchesdisplayedanswer}
The assigned autoformalize agent should translate this physics problem into Lean declarations in the covered file, with theorem and lemma proof bodies written as `by sorry`.
\end{theorem}
\begin{proof}
This is an autoformalization task, not a proof task. Produce statements that can later be proved by the physics prover without weakening the physical model.
\end{proof}

% --- Archon declaration topology begin ---
\section{Declaration topology}
\begin{definition}[Rest Mass Quantity]
  \label{def:physics:phyx-mini-0560:phyxminiproblems-problemphyxmini0560-restmassquantity}
  \lean{PhyXMiniProblems.ProblemPhyXMini0560.RestMassQuantity}
  A nonnegative physical rest mass with mass dimension M
\end{definition}
\begin{proof}
  This is the declared model definition of Rest Mass Quantity.
\end{proof}
\begin{definition}[rest Mass In Kilograms]
  \label{def:physics:phyx-mini-0560:phyxminiproblems-problemphyxmini0560-restmassinkilograms}
  \lean{PhyXMiniProblems.ProblemPhyXMini0560.restMassInKilograms}
  \uses{def:physics:phyx-mini-0560:phyxminiproblems-problemphyxmini0560-restmassquantity}
  Read a physical rest mass in kilograms
\end{definition}
\begin{proof}
  This is the declared model definition of rest Mass In Kilograms.
\end{proof}
\begin{definition}[energy In Joules]
  \label{def:physics:phyx-mini-0560:phyxminiproblems-problemphyxmini0560-energyinjoules}
  \lean{PhyXMiniProblems.ProblemPhyXMini0560.energyInJoules}
  Read a physical energy in joules
\end{definition}
\begin{proof}
  This is the declared model definition of energy In Joules.
\end{proof}
\begin{definition}[speed In Meters Per Second]
  \label{def:physics:phyx-mini-0560:phyxminiproblems-problemphyxmini0560-speedinmeterspersecond}
  \lean{PhyXMiniProblems.ProblemPhyXMini0560.speedInMetersPerSecond}
  Read a nonnegative physical speed in metres per second
\end{definition}
\begin{proof}
  This is the declared model definition of speed In Meters Per Second.
\end{proof}
\begin{definition}[vacuum Light Speed In Meters Per Second]
  \label{def:physics:phyx-mini-0560:phyxminiproblems-problemphyxmini0560-vacuumlightspeedinmeterspersecond}
  \lean{PhyXMiniProblems.ProblemPhyXMini0560.vacuumLightSpeedInMetersPerSecond}
  Physlib's exact dimensionful vacuum speed of light, read in SI units
\end{definition}
\begin{proof}
  This is the declared model definition of vacuum Light Speed In Meters Per Second.
\end{proof}
\begin{definition}[speed Fraction Of Light]
  \label{def:physics:phyx-mini-0560:phyxminiproblems-problemphyxmini0560-speedfractionoflight}
  \lean{PhyXMiniProblems.ProblemPhyXMini0560.speedFractionOfLight}
  \uses{def:physics:phyx-mini-0560:phyxminiproblems-problemphyxmini0560-speedinmeterspersecond, def:physics:phyx-mini-0560:phyxminiproblems-problemphyxmini0560-vacuumlightspeedinmeterspersecond}
  A speed expressed as the dimensionless fraction u/c
\end{definition}
\begin{proof}
  This is the declared model definition of speed Fraction Of Light.
\end{proof}
\begin{definition}[rest Energy Equivalent In Joules]
  \label{def:physics:phyx-mini-0560:phyxminiproblems-problemphyxmini0560-restenergyequivalentinjoules}
  \lean{PhyXMiniProblems.ProblemPhyXMini0560.restEnergyEquivalentInJoules}
  \uses{def:physics:phyx-mini-0560:phyxminiproblems-problemphyxmini0560-restmassquantity, def:physics:phyx-mini-0560:phyxminiproblems-problemphyxmini0560-restmassinkilograms, def:physics:phyx-mini-0560:phyxminiproblems-problemphyxmini0560-vacuumlightspeedinmeterspersecond}
  The rest-energy equivalent m c², expressed in joules
\end{definition}
\begin{proof}
  This is the declared model definition of rest Energy Equivalent In Joules.
\end{proof}
\begin{definition}[Inertial Frame Label]
  \label{def:physics:phyx-mini-0560:phyxminiproblems-problemphyxmini0560-inertialframelabel}
  \lean{PhyXMiniProblems.ProblemPhyXMini0560.InertialFrameLabel}
  The two inertial frames named in the problem
\end{definition}
\begin{proof}
  This is the declared model definition of Inertial Frame Label.
\end{proof}
\begin{definition}[Diagram Axis]
  \label{def:physics:phyx-mini-0560:phyxminiproblems-problemphyxmini0560-diagramaxis}
  \lean{PhyXMiniProblems.ProblemPhyXMini0560.DiagramAxis}
  The two spatial axes shown in the bitmap and used for the boost
\end{definition}
\begin{proof}
  This is the declared model definition of Diagram Axis.
\end{proof}
\begin{definition}[to Fin]
  \label{def:physics:phyx-mini-0560:phyxminiproblems-problemphyxmini0560-diagramaxis-tofin}
  \lean{PhyXMiniProblems.ProblemPhyXMini0560.DiagramAxis.toFin}
  \uses{def:physics:phyx-mini-0560:phyxminiproblems-problemphyxmini0560-diagramaxis}
  Coordinate index corresponding to a labelled diagram axis
\end{definition}
\begin{proof}
  This is the declared model definition of to Fin.
\end{proof}
\begin{definition}[Particle Approach Figure]
  \label{def:physics:phyx-mini-0560:phyxminiproblems-problemphyxmini0560-particleapproachfigure}
  \lean{PhyXMiniProblems.ProblemPhyXMini0560.ParticleApproachFigure}
  \uses{def:physics:phyx-mini-0560:phyxminiproblems-problemphyxmini0560-diagramaxis}
  Qualitative and geometric information transcribed from the primary bitmap. Coordinates here are drawing coordinates only; the bitmap supplies no quantitative position or distance scale
\end{definition}
\begin{proof}
  This is the declared model definition of Particle Approach Figure.
\end{proof}
\begin{definition}[Transverse Boost Setup]
  \label{def:physics:phyx-mini-0560:phyxminiproblems-problemphyxmini0560-transverseboostsetup}
  \lean{PhyXMiniProblems.ProblemPhyXMini0560.TransverseBoostSetup}
  \uses{def:physics:phyx-mini-0560:phyxminiproblems-problemphyxmini0560-restmassquantity, def:physics:phyx-mini-0560:phyxminiproblems-problemphyxmini0560-inertialframelabel, def:physics:phyx-mini-0560:phyxminiproblems-problemphyxmini0560-particleapproachfigure}
  The independent physical quantities of the scenario. In particular, the energy measured in SPrime is an unconstrained dimensionful field here; it is not defined from the requested formula or from an answer choice
\end{definition}
\begin{proof}
  This is the declared model definition of Transverse Boost Setup.
\end{proof}
\begin{definition}[velocity Fraction Component]
  \label{def:physics:phyx-mini-0560:phyxminiproblems-problemphyxmini0560-velocityfractioncomponent}
  \lean{PhyXMiniProblems.ProblemPhyXMini0560.velocityFractionComponent}
  \uses{def:physics:phyx-mini-0560:phyxminiproblems-problemphyxmini0560-diagramaxis, def:physics:phyx-mini-0560:phyxminiproblems-problemphyxmini0560-transverseboostsetup}
  The component u\_axis/c of the particle's velocity in S
\end{definition}
\begin{proof}
  This is the declared model definition of velocity Fraction Component.
\end{proof}
\begin{definition}[energy Momentum Temporal Component]
  \label{def:physics:phyx-mini-0560:phyxminiproblems-problemphyxmini0560-energymomentumtemporalcomponent}
  \lean{PhyXMiniProblems.ProblemPhyXMini0560.energyMomentumTemporalComponent}
  \uses{def:physics:phyx-mini-0560:phyxminiproblems-problemphyxmini0560-inertialframelabel, def:physics:phyx-mini-0560:phyxminiproblems-problemphyxmini0560-transverseboostsetup}
  Temporal energy coordinate of the energy-momentum readout, in joules
\end{definition}
\begin{proof}
  This is the declared model definition of energy Momentum Temporal Component.
\end{proof}
\begin{definition}[energy Momentum Spatial Component]
  \label{def:physics:phyx-mini-0560:phyxminiproblems-problemphyxmini0560-energymomentumspatialcomponent}
  \lean{PhyXMiniProblems.ProblemPhyXMini0560.energyMomentumSpatialComponent}
  \uses{def:physics:phyx-mini-0560:phyxminiproblems-problemphyxmini0560-inertialframelabel, def:physics:phyx-mini-0560:phyxminiproblems-problemphyxmini0560-diagramaxis, def:physics:phyx-mini-0560:phyxminiproblems-problemphyxmini0560-transverseboostsetup}
  Spatial p\_axis c coordinate of the energy-momentum readout, in joules
\end{definition}
\begin{proof}
  This is the declared model definition of energy Momentum Spatial Component.
\end{proof}
\begin{definition}[Matches Problem And Primary Figure]
  \label{def:physics:phyx-mini-0560:phyxminiproblems-problemphyxmini0560-matchesproblemandprimaryfigure}
  \lean{PhyXMiniProblems.ProblemPhyXMini0560.MatchesProblemAndPrimaryFigure}
  \uses{def:physics:phyx-mini-0560:phyxminiproblems-problemphyxmini0560-speedfractionoflight, def:physics:phyx-mini-0560:phyxminiproblems-problemphyxmini0560-diagramaxis, def:physics:phyx-mini-0560:phyxminiproblems-problemphyxmini0560-transverseboostsetup, def:physics:phyx-mini-0560:phyxminiproblems-problemphyxmini0560-velocityfractioncomponent}
  Facts read from the primary raster and the prose. The arrow starts at the particle in the first quadrant and points inward along the same ray. The component equations encode an inward velocity making alpha with the positive y-axis. No energy in SPrime is fixed here
\end{definition}
\begin{proof}
  This is the declared model definition of Matches Problem And Primary Figure.
\end{proof}
\begin{definition}[Has Physical Input Parameters]
  \label{def:physics:phyx-mini-0560:phyxminiproblems-problemphyxmini0560-hasphysicalinputparameters}
  \lean{PhyXMiniProblems.ProblemPhyXMini0560.HasPhysicalInputParameters}
  \uses{def:physics:phyx-mini-0560:phyxminiproblems-problemphyxmini0560-restmassinkilograms, def:physics:phyx-mini-0560:phyxminiproblems-problemphyxmini0560-energyinjoules, def:physics:phyx-mini-0560:phyxminiproblems-problemphyxmini0560-vacuumlightspeedinmeterspersecond, def:physics:phyx-mini-0560:phyxminiproblems-problemphyxmini0560-speedfractionoflight, def:physics:phyx-mini-0560:phyxminiproblems-problemphyxmini0560-transverseboostsetup}
  Positivity, angle range, and subluminal input conditions. They select the physical branch but contain no equation for the requested transformed energy
\end{definition}
\begin{proof}
  This is the declared model definition of Has Physical Input Parameters.
\end{proof}
\begin{definition}[Observer Makes Particle Move Along YPrime]
  \label{def:physics:phyx-mini-0560:phyxminiproblems-problemphyxmini0560-observermakesparticlemovealongyprime}
  \lean{PhyXMiniProblems.ProblemPhyXMini0560.ObserverMakesParticleMoveAlongYPrime}
  \uses{def:physics:phyx-mini-0560:phyxminiproblems-problemphyxmini0560-transverseboostsetup, def:physics:phyx-mini-0560:phyxminiproblems-problemphyxmini0560-velocityfractioncomponent, def:physics:phyx-mini-0560:phyxminiproblems-problemphyxmini0560-energymomentumspatialcomponent}
  The defining condition on the observer in SPrime. Its boost is along the x direction and matches the particle's transverse velocity fraction, while the transformed particle has zero xPrime momentum and hence travels along the yPrime axis. This specifies the observer, not the requested energy
\end{definition}
\begin{proof}
  This is the declared model definition of Observer Makes Particle Move Along YPrime.
\end{proof}
\begin{definition}[Satisfies Special Relativistic Particle Kinematics]
  \label{def:physics:phyx-mini-0560:phyxminiproblems-problemphyxmini0560-satisfiesspecialrelativisticparticlekinematics}
  \lean{PhyXMiniProblems.ProblemPhyXMini0560.SatisfiesSpecialRelativisticParticleKinematics}
  \uses{def:physics:phyx-mini-0560:phyxminiproblems-problemphyxmini0560-energyinjoules, def:physics:phyx-mini-0560:phyxminiproblems-problemphyxmini0560-restenergyequivalentinjoules, def:physics:phyx-mini-0560:phyxminiproblems-problemphyxmini0560-inertialframelabel, def:physics:phyx-mini-0560:phyxminiproblems-problemphyxmini0560-diagramaxis, def:physics:phyx-mini-0560:phyxminiproblems-problemphyxmini0560-transverseboostsetup, def:physics:phyx-mini-0560:phyxminiproblems-problemphyxmini0560-velocityfractioncomponent, def:physics:phyx-mini-0560:phyxminiproblems-problemphyxmini0560-energymomentumtemporalcomponent, def:physics:phyx-mini-0560:phyxminiproblems-problemphyxmini0560-energymomentumspatialcomponent}
  The four-vector readouts use energy coordinates (E, p\_x c, p\_y c). Accordingly: the temporal coordinate is the measured relativistic energy; in S, p\_i c = E (u\_i/c); the change from S to SPrime is Physlib's Lorentz boost; the invariant mass shell is E² - |p c|² = (m c²)². These are general governing laws. None states the target square-root formula for the energy in SPrime
\end{definition}
\begin{proof}
  This is the declared model definition of Satisfies Special Relativistic Particle Kinematics.
\end{proof}
\begin{lemma}[energy after transverse boost]
  \label{lem:physics:phyx-mini-0560:phyxminiproblems-problemphyxmini0560-energy-after-transverse-boost}
  \lean{PhyXMiniProblems.ProblemPhyXMini0560.energy_after_transverse_boost}
  \uses{def:physics:phyx-mini-0560:phyxminiproblems-problemphyxmini0560-energyinjoules, def:physics:phyx-mini-0560:phyxminiproblems-problemphyxmini0560-transverseboostsetup, def:physics:phyx-mini-0560:phyxminiproblems-problemphyxmini0560-velocityfractioncomponent, def:physics:phyx-mini-0560:phyxminiproblems-problemphyxmini0560-hasphysicalinputparameters, def:physics:phyx-mini-0560:phyxminiproblems-problemphyxmini0560-observermakesparticlemovealongyprime, def:physics:phyx-mini-0560:phyxminiproblems-problemphyxmini0560-satisfiesspecialrelativisticparticlekinematics}
  Before inserting the figure's trigonometric component, the transverse boost changes the energy by sqrt (1 - (u\_x/c)\textasciicircum{}2)
\end{lemma}
\begin{proof}
  The conclusion follows from the governing relations and figure readouts named in the statement of energy after transverse boost.
\end{proof}
\begin{definition}[Answer Choice]
  \label{def:physics:phyx-mini-0560:phyxminiproblems-problemphyxmini0560-answerchoice}
  \lean{PhyXMiniProblems.ProblemPhyXMini0560.AnswerChoice}
  Labels of the four displayed multiple-choice answers
\end{definition}
\begin{proof}
  This is the declared model definition of Answer Choice.
\end{proof}
\begin{definition}[displayed Energy Multiplier]
  \label{def:physics:phyx-mini-0560:phyxminiproblems-problemphyxmini0560-displayedenergymultiplier}
  \lean{PhyXMiniProblems.ProblemPhyXMini0560.displayedEnergyMultiplier}
  \uses{def:physics:phyx-mini-0560:phyxminiproblems-problemphyxmini0560-speedfractionoflight, def:physics:phyx-mini-0560:phyxminiproblems-problemphyxmini0560-transverseboostsetup, def:physics:phyx-mini-0560:phyxminiproblems-problemphyxmini0560-answerchoice}
  The dimensionless multiplier of E printed beside each answer label
\end{definition}
\begin{proof}
  This is the declared model definition of displayed Energy Multiplier.
\end{proof}
\begin{definition}[recorded Dataset Answer]
  \label{def:physics:phyx-mini-0560:phyxminiproblems-problemphyxmini0560-recordeddatasetanswer}
  \lean{PhyXMiniProblems.ProblemPhyXMini0560.recordedDatasetAnswer}
  \uses{def:physics:phyx-mini-0560:phyxminiproblems-problemphyxmini0560-answerchoice}
  The answer label recorded in the source dataset; it is not a premise
\end{definition}
\begin{proof}
  This is the declared model definition of recorded Dataset Answer.
\end{proof}
\begin{definition}[Matches Displayed Answer]
  \label{def:physics:phyx-mini-0560:phyxminiproblems-problemphyxmini0560-matchesdisplayedanswer}
  \lean{PhyXMiniProblems.ProblemPhyXMini0560.MatchesDisplayedAnswer}
  \uses{def:physics:phyx-mini-0560:phyxminiproblems-problemphyxmini0560-energyinjoules, def:physics:phyx-mini-0560:phyxminiproblems-problemphyxmini0560-transverseboostsetup, def:physics:phyx-mini-0560:phyxminiproblems-problemphyxmini0560-answerchoice, def:physics:phyx-mini-0560:phyxminiproblems-problemphyxmini0560-displayedenergymultiplier}
  A displayed multiplier agrees with the independently modeled energy
\end{definition}
\begin{proof}
  This is the declared model definition of Matches Displayed Answer.
\end{proof}
\begin{theorem}[problem phyx mini 0560]
  \label{thm:physics:phyx-mini-0560:phyxminiproblems-problemphyxmini0560-problem-phyx-mini-0560}
  \lean{PhyXMiniProblems.ProblemPhyXMini0560.problem_phyx_mini_0560}
  \uses{def:physics:phyx-mini-0560:phyxminiproblems-problemphyxmini0560-energyinjoules, def:physics:phyx-mini-0560:phyxminiproblems-problemphyxmini0560-speedfractionoflight, def:physics:phyx-mini-0560:phyxminiproblems-problemphyxmini0560-transverseboostsetup, def:physics:phyx-mini-0560:phyxminiproblems-problemphyxmini0560-matchesproblemandprimaryfigure, def:physics:phyx-mini-0560:phyxminiproblems-problemphyxmini0560-hasphysicalinputparameters, def:physics:phyx-mini-0560:phyxminiproblems-problemphyxmini0560-observermakesparticlemovealongyprime, def:physics:phyx-mini-0560:phyxminiproblems-problemphyxmini0560-satisfiesspecialrelativisticparticlekinematics}
  With u\_x/c = -(u/c) sin alpha, the energy seen in SPrime is EPrime = E sqrt (1 - (u²/c²) sin² alpha), which is the expression printed as answer D. This formalizes
\end{theorem}
\begin{proof}
  The conclusion follows from the governing relations and figure readouts named in the statement of problem phyx mini 0560.
\end{proof}
% --- Archon declaration topology end ---
% --- Archon physics formalization source end ---
